% Generated draft: requires mathematical and source review.
\documentclass[11pt]{article}
\usepackage[margin=25mm]{geometry}
\usepackage{fontspec}
\setmainfont{Latin Modern Roman}
\usepackage{amsmath,amssymb,mathtools,unicode-math}
\setmathfont{Latin Modern Math}
\usepackage{enumitem,needspace,xurl,hyperref}
\setlength{\parindent}{0pt}
\setlength{\parskip}{5pt}
\setlength{\emergencystretch}{3em}
\newcommand{\N}{\mathbb{N}}
\newcommand{\Z}{\mathbb{Z}}
\newcommand{\Q}{\mathbb{Q}}
\newcommand{\R}{\mathbb{R}}
\newcommand{\C}{\mathbb{C}}
\newcommand{\F}{\mathbb{F}}
\newcommand{\A}{\mathbb{A}}
\newcommand{\PP}{\mathbb{P}}
\newcommand{\E}{\mathbb{E}}
\newcommand{\eps}{\varepsilon}
\newcommand{\dd}{\,\mathrm{d}}
\begin{document}
\section*{UnsolvedMath ID: 700263}
\textbf{Source problem code:} catalogue entry
\subsection{Definitions and mathematical statement}
Let \(R\) be a commutative Noetherian ring with identity, and let \(\Lambda\) be an \(R\)-algebra that is finitely generated as an \(R\)-module. Write \(\operatorname{Max}(R)\) for the set of maximal ideals of \(R\), and let \(\Lambda^{(n)}\) denote the direct sum of \(n\) copies of \(\Lambda\). For an \(R\)-module \(M\) and \(\mathfrak m\in\operatorname{Max}(R)\), put
\[
M_{\mathfrak m}=(R\setminus\mathfrak m)^{-1}M.
\]
The source uses \(\operatorname{Mod}\Lambda\) without specifying whether its objects are left or right modules; a consistent choice of side is required throughout.

Let \(T\) be the full subcategory of \(\operatorname{Mod}\Lambda\) consisting of modules \(U\) for which there exist a finite nonnegative integer \(n\), a module \(P\), and an exact sequence of \(\Lambda\)-modules
\[
0\longrightarrow P\longrightarrow\Lambda^{(n)}\longrightarrow U\longrightarrow0
\]
such that
\[
P_{\mathfrak m}\cong\Lambda_{\mathfrak m}^{(n)}
\qquad\text{for every }\mathfrak m\in\operatorname{Max}(R).
\]
Here the displayed local isomorphisms are isomorphisms of \(\Lambda_{\mathfrak m}\)-modules. Define \(\mathbb G(T)\) to be the abelian group generated by the isomorphism classes \([U]\) of objects of \(T\), subject to the relations
\[
[B]=[A]+[C]
\]
for every short exact sequence \(0\to A\to B\to C\to0\) of \(\Lambda\)-modules whose three terms belong to \(T\).

\begin{enumerate}
\item[(a)] Describe \(\mathbb G(T)\). In particular, characterize, for \(U,V\in T\), when \([U]=[V]\) in \(\mathbb G(T)\).
\item[(b)] The source proposes the following conjecture. Suppose
\[
d=\dim\bigl(\operatorname{Max}(R)\bigr)<\infty.
\]
For every \(U,V\in T\), every finite integer \(n>d\), and every pair of epimorphisms of \(\Lambda\)-modules
\[
\varphi:\Lambda^{(n)}\twoheadrightarrow U,
\qquad
\psi:\Lambda^{(n)}\twoheadrightarrow V,
\]
is it true that
\[
[U]=[V]\text{ in }\mathbb G(T)
\quad\Longrightarrow\quad
\ker\varphi=\ker\psi?
\]
\end{enumerate}
The conclusion in (b) reproduces the source's literal equality of submodules of the common domain \(\Lambda^{(n)}\). The supplied record does not clarify whether equality, abstract isomorphism, or equivalence under an automorphism of the domain was intended; these conclusions must not be identified. The source also does not specify the dimension convention in \(\dim(\operatorname{Max}(R))\), which is retained here without replacing it by \(\dim R\). The epimorphisms in (b) are quantified as written; the source does not explicitly require them to be the epimorphisms in the presentations defining \(T\).
\subsection{Short English statement}
For a module-finite algebra over a commutative Noetherian ring, consider modules admitting a finite free presentation whose kernel becomes free of the same rank at every maximal ideal. Describe their Grothendieck group and equality of its classes. The source further conjectures that equal classes force literal equality of kernels of any two epimorphisms from the same free module of rank greater than the stated dimension of the maximal spectrum. The intended kernel comparison, dimension convention, and module handedness require clarification.
\subsection{Sources}
\begin{itemize}
\item Supplied local UnsolvedMath registry record: permanent ID \(700263\), exact source problem code \texttt{catalogue entry}, title \emph{Kourovka Notebook Problem 4.56}; upstream revision \texttt{d371ee917ac9647fbb06cd6fc529deec6f3dda82}.
\item \emph{The Kourovka Notebook (Unsolved Problems in Group Theory)}, Problem 4.56, attributed to K. W. Roggenkamp, proposed in 1973 and originating in the fourth issue (1973). The supplied background identifies the twenty-first-issue edition (2026), arXiv:1401.0300v48: \url{https://arxiv.org/abs/1401.0300v48}. The supplied collection description names E. I. Khukhro and V. D. Mazurov as editors of the Issues 1--21 collection (1965--2026).
\item The registry labels this record \emph{open}. This is a supplied status label, not an independently verified assessment of the literature. The cited source has not been independently consulted for this draft.
\end{itemize}
\end{document}
